\section{Выводы}

В ходе выполнения лабораторной работы изучены и использованы штатные средства работы с
потоками ОС Linux: \texttt{pthread\_create}, \texttt{pthread\_join}, мьютексы
\texttt{pthread\_mutex\_*} и барьеры \texttt{pthread\_barrier\_t}, а также измерение времени
через \texttt{clock\_gettime(CLOCK\_MONOTONIC)}.

Разработана многопоточная программа варианта 11: K-кратное наложение медианного фильтра на
матрицу целых чисел с задаваемым пользовательским окном. Матрица разбивается на
непересекающиеся блоки строк, которые захватываются потоками из общей очереди под мьютексом;
границы итераций синхронизируются барьером, необходимым для корректного переключения двойного
буфера \texttt{a}/\texttt{b}. Лимит \texttt{--threads P} ограничивает общее число вычислительных
потоков вместе с главным, при $P=1$ программа работает полностью последовательно.

Корректность подтверждена сверкой с последовательным эталоном (\texttt{--check}) при всех
значениях P; число потоков продемонстрировано штатными средствами ОС (\texttt{ps -eLf},
\texttt{top -H}, \texttt{/proc/<pid>/status}). Экспериментально получены ускорение
$S(P)=T_1/T_P$ и эффективность $E(P)=100\,S(P)/P$: на 8 потоках достигается пятикратное
ускорение, при дальнейшем росте P выгоды сокращаются из-за накладных расходов на синхронизацию
и уменьшения доли работы на поток. Сформирован полноценный отчёт о ходе и результатах работы.
\pagebreak
